OpenFisca-France (législation des revenus 2024) ; DGFiP, IRCOM revenus 2024 ; calculs de l'auteur (synthese/modele\_purs.py, architectures\_ircsg.py).\par Note : impôt individuel sur le revenu après cotisations, même barème pour tous, même recette qu'aujourd'hui (220 Md\euro{}). Taux marginal par adulte en ligne droite entre les seuils. Barème neutre (le plus proche de ce que chaque foyer paie aujourd'hui) : 2 \% au premier euro, 14 \% à 20 000 \euro{}, 23 \% à 30 000 \euro{}, 50 \% à 250 000 \euro{}. Proposition : franchise de 4 000 \euro{} par adulte sur les revenus d'activité, 6 000 \euro{} par enfant, pensions exonérées jusqu'à 12 500 \euro{} (puis de moins en moins jusqu'à 25 000 \euro{}) ; 14 \% au premier euro imposé, 14 \% à 20 000 \euro{}, 30,3 \% à 45 000 \euro{}, 50 \% à 1000 000 \euro{}. Cas-types sans niche : IR + CSG + CRDS (et Casa pour le retraité) en \% du revenu brut ; taux marginal sur 1 000 \euro{} de plus, lissé sur 2 000 \euro{}.
